\section*{Appendix}

\subsection{Proof of Single-Sideband Modulation}
\label{subsec:proof_of_ssb}
Assume the high-frequency carrier signal is $c(t) = cos(2\pi f_c t)$ and the noise signal is $n(t)$, the DSB-AM and LSB-AM can be represented by Equation \ref{equ:DSB} and \ref{equ:LSB}:
\begin{equation}
   n_D(t) = \sqrt{2}n(t) cos(2\pi f_c t)
    \label{equ:DSB}
\end{equation}
\begin{equation}
    n_L(t) = n(t) cos(2\pi f_c t) + \widehat{n}(t) sin(2\pi f_c t)
    \label{equ:LSB}
 \end{equation}
where $f_c$ is the carrier frequency, and $\hat{n}(t)$ is the Hilbert transform of $n(t)$. The coefficient $\sqrt{2}$ in Equation \ref{equ:DSB} is to ensure the energy of the two modulated signals are same.
 
Similar to microphones, the nonlinearity in the ultrasonic transmitter also generates high order components. Because energy decays significantly as the order goes up, we only consider the quadratic term here. With DSB-AM and LSB-AM modulation, the quadratic terms of the modulated signal can be represented by Equation \ref{equ:DSB_qua} and \ref{equ:LSB_qua}. As shown in these two equations, the human audible low-frequency components for DSB-AM is $n^2(t)$, while for LSB-AM it is $\frac{1}{2}(n^2(t) + \widehat{n}^2(t))$. Because the Hilbert transform imparts a phase shift of $\pm \frac{\pi}{2}$ to each frequency components, the amplitude of each frequency components in the former is $\sqrt{2}$ times bigger than the latter. That is, the former has twice the energy of the latter. 
\begin{equation}
    \begin{aligned}
     n^2_{D}(t) = &n^2(t)(1+cos(4\pi f_c t))\xRightarrow[Filter]{Lowpass} n^2(t)
    \end{aligned}
    \label{equ:DSB_qua}
\end{equation}
\begin{equation}
    \begin{aligned}
        n^2_{L}(t) \xRightarrow[Filter]{Lowpass}0.5(n^2(t) + \widehat{n}^2(t))
    \end{aligned}
    \label{equ:LSB_qua}
\end{equation}
